% Generated draft: requires mathematical and source review.
\documentclass[11pt]{article}
\usepackage[margin=25mm]{geometry}
\usepackage{fontspec}
\setmainfont{Latin Modern Roman}
\usepackage{amsmath,amssymb,mathtools,unicode-math}
\setmathfont{Latin Modern Math}
\usepackage{enumitem,needspace,xurl,hyperref}
\setlength{\parindent}{0pt}
\setlength{\parskip}{5pt}
\setlength{\emergencystretch}{3em}
\newcommand{\N}{\mathbb{N}}
\newcommand{\Z}{\mathbb{Z}}
\newcommand{\Q}{\mathbb{Q}}
\newcommand{\R}{\mathbb{R}}
\newcommand{\C}{\mathbb{C}}
\newcommand{\F}{\mathbb{F}}
\newcommand{\A}{\mathbb{A}}
\newcommand{\PP}{\mathbb{P}}
\newcommand{\E}{\mathbb{E}}
\newcommand{\eps}{\varepsilon}
\newcommand{\dd}{\,\mathrm{d}}
\begin{document}
\section*{UnsolvedMath ID: 30004611}
\textbf{Source problem code:} catalogue entry
\subsection{Definitions and mathematical statement}
The supplied question concerns a hyperplane arrangement \(\mathcal A\) with a distinguished hyperplane \(H\in\mathcal A\). Its deletion is \(\mathcal A'=\mathcal A\setminus\{H\}\), and its restriction is the arrangement of distinct hyperplanes \(\mathcal A''=\{H\cap K:K\in\mathcal A'\}\) in \(H\), whenever these intersections are hyperplanes in \(H\). The source does not specify the ground field, centrality, ambient space, or sheaf convention.

For clarity, the standard central-arrangement definitions can be given conditionally. Suppose \(\mathcal A\) is a finite arrangement of distinct linear hyperplanes in a finite-dimensional vector space \(V\) over a field \(k\). Put \(S=\operatorname{Sym}_k(V^*)\), and choose a nonzero linear defining equation \(\alpha_K\) for each \(K\in\mathcal A\). The module of logarithmic derivations is
\[
D(\mathcal A)=\{\theta\in\operatorname{Der}_k(S):\theta(\alpha_K)\in\alpha_KS\text{ for every }K\in\mathcal A\}.
\]
Define \(D(\mathcal A')\) by the same rule for the deletion, and define \(D(\mathcal A'')\) over \(S_H=\operatorname{Sym}_k(H^*)\). The Ziegler multiplicity on \(\mathcal A''\) is
\[
m_H(X)=\#\{K\in\mathcal A':H\cap K=X\},\qquad X\in\mathcal A''.
\]
If \(\beta_X\in H^*\) defines \(X\), the corresponding multiarrangement module is
\[
D(\mathcal A'',m_H)=\{\eta\in\operatorname{Der}_k(S_H):\eta(\beta_X)\in\beta_X^{m_H(X)}S_H\text{ for every }X\in\mathcal A''\}.
\]
On the affine spaces \(V=\operatorname{Spec}S\) and \(H=\operatorname{Spec}S_H\), these modules have associated quasi-coherent sheaves. In this conditional convention, the source notation \(D,D',D''\) corresponds to the sheaves associated with \(D(\mathcal A),D(\mathcal A'),D(\mathcal A'')\), respectively. Its notation \((D_H,m_H)\) appears to refer to the logarithmic derivation sheaf for the Ziegler multiarrangement; this interpretation is not explicitly defined in the supplied record.

\textbf{Source question.} Are there short exact sequences relating the logarithmic derivation sheaves of an arrangement, its deletion, its restriction, and its Ziegler restriction? A short exact sequence means a sequence of sheaf morphisms
\[
0\longrightarrow\mathcal F\xrightarrow{u}\mathcal G\xrightarrow{v}\mathcal K\longrightarrow0
\]
with \(u\) injective, \(\operatorname{im}u=\ker v\), and \(v\) surjective. The source does not specify the proposed terms, morphisms, twists, hypotheses, or whether several sequences are intended. Restriction sheaves live on \(H\); a sequence on the ambient space would require an explicit convention such as direct image under the inclusion \(H\hookrightarrow V\). The conditional definitions above explain the terminology but do not establish that the source intended the central affine setting. Consequently, the supplied question remains a research target with unresolved formulation details, rather than a fully specified exactness assertion.
\subsection{Short English statement}
The problem asks for short exact sequences connecting logarithmic derivation sheaves before and after deleting or restricting a hyperplane, including the multiplicities of Ziegler restriction. The source leaves the setting, maps, and hypotheses unspecified. Its supplied status is partially solved; this draft does not independently verify that assessment.
\subsection{Sources}
\begin{itemize}
\item Supplied local UnsolvedMath registry record, permanent ID \(30004611\), title \emph{Exact Sequences for Logarithmic Derivation Sheaves}, source problem code \texttt{catalogue entry}, in the Oberwolfach Reports Open Problems collection. The record gives proposed year 2021 and names no individual proposer.
\item Supplied report citation: \emph{Logarithmic Vector Fields and Freeness of Divisors and Arrangements: New perspectives and applications} (2022), Oberwolfach Reports, pp. 225--300, DOI \texttt{10.4171/owr/2021/5}. \url{https://doi.org/10.4171/owr/2021/5}. The supplied background attributes the question to workshop open-problem discussions documented in Oberwolfach Reports.
\item Supplied literature source: \emph{Local cohomology of logarithmic forms}. \url{https://www.math.purdue.edu/~walther/research/logforms.pdf}.
\item Supplied literature source: \emph{Oberwolfach Report 5/2021}. \url{https://ems.press/content/serial-article-files/46883}.
\item Literature assessment checked at \texttt{2026-08-22}: the supplied dated curation labels the problem partially solved and reports established short exact sequences for logarithmic derivations and forms in important special cases, such as non-Boolean generic arrangements, while describing the general deletion, restriction, and Ziegler-restriction question as unresolved. This is supplied dated curation derived from the OWR record, not an independently verified literature assessment; its status and cited claims require independent verification.
\item The registry labels the statement \texttt{corrected\_verified} and reports a manual semantic, typographic, and LaTeX review asserting self-containedness and preservation of the target. This is supplied editorial metadata, not independent verification of the formulation or mathematical claims.
\end{itemize}
\end{document}
